\chapter{Graph-Structured Lambda Theories}
\label{ch:gslt}

\section{The lineage: structured operational semantics}

The presentation style we depend on is Plotkin's~\cite{Plotkin1981,Plotkin}. Before
structured operational semantics, giving the meaning of a programming construct meant
either a denotational assignment into a domain or an informal appeal to an abstract
machine. Plotkin's proposal was that one could give the behaviour of a language
\emph{syntactically}, by inference rules over the term structure, with the rules for a
compound term given in terms of the rules for its parts. The presentation was finite,
it was checkable, and --- this is the part that matters here --- it was
\emph{structural}: the shape of the rules followed the shape of the grammar.

Milner's \emph{Functions as Processes}~\cite{Milner1992} adopts this style and sets the
standard by which such presentations should be judged. The paper gives the
$\pi$-calculus by a grammar, a structural congruence, and a small family of reduction
rules, and then encodes the $\lambda$-calculus into it. What makes the presentation
exemplary is not merely that it is correct but that it is \emph{complete in the
relevant sense}: the grammar, the congruence, and the rules together are the whole of
the definition. Nothing is left to the reader's operational imagination. One can
mechanically check a claimed reduction, and one can mechanically check that an encoding
respects reduction, because there is nothing to the language beyond what has been
written down.

Cardelli and Gordon's presentation of the mobile ambient
calculus~\cite{CardelliGordon2000} adopts precisely this style, and it is instructive
to see how little has to change. The ambient calculus is a very different theory ---
its reductions restructure a spatial hierarchy rather than substitute a value --- and
yet the presentation is recognisably the same document: sorts, term formers, a
structural congruence, a handful of reduction rules, and congruence rules closing
reduction under the term formers. The style is not an artifact of process calculi. It
is a general format for presenting a language.

The observation on which this document rests is that once a presentation is
\emph{that} regular, it is data. A grammar is a signature. A structural congruence is a
set of equations. Reduction rules and congruence rules are a set of rewrites. These are
not analogies; they are the same information under different typography. A GSLT is what
you get when you stop treating an SOS presentation as a document to be read by a human
implementer and start treating it as a structure to be consumed by a machine.

\section{The definition}

\begin{definition}[Graph-structured lambda theory]
\label{ob:def:gslt}
A \textbf{graph-structured lambda theory} is a triple $G = (\Sigma, E, R)$ consisting of
\begin{itemize}[leftmargin=2em]
  \item a signature $\Sigma$, which is in general a \emph{lambda theory} --- so that
    terms may contain variables;
  \item a set $E$ of equations, imposing a structural congruence $\equiv$; and
  \item a set $R$ of rewrite rules.
\end{itemize}
\end{definition}

\begin{remark}[A bridge between two spellings]
\label{ob:rmk:notation-bridge}
This part writes a theory as $G = (\Sigma, E, R)$ and reserves boldface
for the categories --- $\catGSLT$, $\catiGSLT$, $\catciGSLT$ --- because
that is the notation in which the constructions were first worked out and
in which the reader will find them in the reference literature.
Later parts of the book write the same data as
$\GSLT = (\terms, \eqs, \rules)$, a theory being denoted by a single
calligraphic letter because the parts that use it care about which
\emph{theory} they are in far more often than they care about which
component of it they are looking at.
Nothing turns on the difference.
Where a later chapter says ``the GSLT $\GSLT$'' it means an object of
$\catGSLT$, and where this part says ``the signature $\Sigma$'' a later
chapter would say ``the grammar $\terms$''.
\end{remark}

The definition is deliberately thin in one direction and deliberately generous in
another, and both deserve comment.

It is thin in that a GSLT is any rewrite theory whatever. Nothing in
Definition~\ref{ob:def:gslt} says the theory is about computation, or that it has processes,
or that anything interacts with anything. The strengthenings of
Chapters~\ref{ch:igslt} and~\ref{ch:cigslt} add exactly the structure later
constructions require, and the discipline of the present document is to add it only when
a construction needs it.\footnote{We say \emph{strengthening} rather than
\emph{enrichment} throughout, reserving the latter for its technical sense, which
Section~\ref{ob:subsec:name} argues is not the right technique here.}

It is generous in that $\Sigma$ is a lambda theory rather than a bare algebraic
signature. That is the substantive choice in the definition, and the next subsection is
about why it was made.

\section{Reading the name}
\label{ob:subsec:name}

The two halves of ``graph-structured lambda theory'' are doing two separable jobs. The
\emph{lambda theory} captures the term language --- specifically, a term language
\emph{with variables}. The adjective \emph{graph-structured} captures the rewrite rules.
It is worth taking these in turn, and worth recording the route not taken, because the
route not taken is the one a reader familiar with categorical algebra will expect.

\section{Why lambda theories and not Lawvere theories}

The natural first thought is that a language presentation is an algebraic theory, and
that the right categorical home for an algebraic theory is a Lawvere
theory~\cite{Lawvere1963}: a category with finite products whose objects are the powers
of a generic object, in which an $n$-ary operation is a morphism $n \to 1$. This is a
beautiful and well-understood setting, and it is not adequate here.

The reason is variables. A Lawvere theory does not generate a term language with
variables; it generates a term language in which variables have been \emph{absorbed into
arities}. A term in $n$ variables is not a term containing variables, it is a morphism
out of $n$. That encoding is faithful for first-order algebra and it breaks down exactly
where we need it not to: one cannot abstract, one cannot have a variable that ranges over
functions, and one cannot write a binder. Every calculus this document treats has
binders --- $\lambda x.M$, $\mathrm{for}(x \leftarrow n)P$, $\mathrm{new}(x, P)$ --- and
a formalism that cannot express a binder cannot present them.

Lambda theories do generate term languages with variables. And in the limit --- that is,
in how the variables get used, which is to say abstraction and application --- lambda
theories are cartesian closed categories. This is the content of the
Curry--Howard--Lambek correspondence~\cite{LambekScott1986}, and it is why the
classifying form of Section~\ref{ob:subsec:classifying} carries a chosen cartesian closed
structure. The CCC structure is not an extra assumption bolted on for the convenience of
Chapter~\ref{ch:hypercube}; it is what a lambda theory becomes when one asks what its
variables are for.

\section{Why adjoined and not enriched}

The second natural thought concerns the rewrites. Rewriting looks like $2$-dimensional
structure --- between two terms there is not merely the question of whether they are
equal but a collection of ways of getting from one to the other --- and the standard
technique for such structure is \emph{enrichment}: take hom-objects in the category of
graphs rather than in $\mathbf{Set}$. Graph-enriched Lawvere theories were the initial
approach, and they too are not adequate here.

The difficulty is that enrichment puts the rewrite structure in the wrong place. Under
enrichment the graph of rewrites lives in the hom, where it is available to whatever
respects composition and to nothing else. But we want to \emph{quantify over} rewrites:
the history monad of Chapter~\ref{ch:histmonad} makes rewrite events into terms, and the
type-generation of Chapter~\ref{ch:hypercube} reads a rewrite rule together with a
choice of position inside its left-hand side. Neither operation is available to a
construction that can only see hom-objects.

The approach we take instead is to \emph{adjoin} the graph structure to the theory,
as a Lawvere theory of graphs. One says what the sources and the targets are, and ensures
that the source and target maps enjoy the minimal requisite equations. The rewrite
structure thereby becomes ordinary algebraic structure inside the theory, presented by
the same mechanism that presents everything else --- rather than a change of base for
enrichment, sitting outside the theory and constraining it from without. The framing of
operational semantics in this algebraic style follows Stay and
Meredith~\cite{OSLF} and Baez and Williams~\cite{BaezWilliams2020}.

\begin{remark}[The name, in a nutshell]
\label{ob:rem:nutshell}
\emph{Lambda theory} captures term languages with variables. \emph{Graph-structured}
captures the rewrite rules. A GSLT is a term language with variables, together with an
adjoined graph of rewrites over it.
\end{remark}

\begin{remark}[Why this matters to the reader who wants to ship]
The choice is visible in the specification format. Because the graph structure is
adjoined rather than enriched, the \texttt{rewrites} block of a specification is not a
different \emph{kind} of declaration from the \texttt{terms} block --- it is more
algebraic structure over the same signature. That is why the ladder of
Section~\ref{ob:subsec:ladder} is a ladder rather than three unrelated facilities, and why a
tool that can read the first two blocks can read the third.
\end{remark}

\section{Presentation and classifying theory}
\label{ob:subsec:classifying}

Definition~\ref{ob:def:gslt} is syntactic: it is what a developer writes and what a machine
accepts. Several constructions below are more naturally stated semantically, and
Section~\ref{ob:subsec:name} has already said what the semantic form is. We record it
explicitly.

The \emph{classifying theory} of a presentation is a category $\mathsf{T}$ with finite
limits and a chosen cartesian closed structure --- the cartesian closure being what the
lambda theory $\Sigma$ becomes in the limit --- equipped with its subobject fibration
$\pi : \mathrm{Sub}(\mathsf{T}) \to \mathsf{T}$, a distinguished object $\Pr$ of
programs, and a distinguished subobject
\[
  (\rsq) \rightarrowtail \Pr \times \Pr
\]
of one-step reduction, together with entailments generating the base rewrites and their
closure under the term formers. This is the presentation-independent form in which the
type-theoretic construction of Chapter~\ref{ch:hypercube} takes its input.

The subobject $(\rsq)$ is the adjoined graph of Section~\ref{ob:subsec:name} seen
semantically: its source and target maps are the two projections, and the entailments
generating the base rewrites and closing them under the term formers are the minimal
equations those maps must satisfy. So the two presentations are not two objects to be
reconciled; they are the syntactic and semantic faces of the same construction, and the
reason the classifying theory has exactly the structure it has is that $\Sigma$ is a
lambda theory and $R$ is adjoined as a graph.

We take the presentation as primary throughout, for the practical reason that it is what
one can type, and use the classifying theory when a construction is more naturally stated
in the internal language.

\section{The ladder}
\label{ob:subsec:ladder}

Definition~\ref{ob:def:gslt} has three components, and the reader can switch two of them
off. Doing so produces a ladder of three rungs, each of which is a familiar and
independently useful thing. We climb it because it is the most efficient route we know
from ``I already know what this is'' to ``I did not know you could do that.''

\begin{center}
\begin{tikzpicture}[
  rung/.style={draw,rounded corners=2pt,minimum height=9mm,inner xsep=8pt,align=left},
  every node/.style={font=\small}
]
  \node[rung] (r1) at (0,0)
    {\textbf{terms only} \;$\longrightarrow$\; algebraic data types};
  \node[rung] (r2) at (0,-1.35)
    {\textbf{$+$ equations} \;$\longrightarrow$\; finitely presentable universal algebra};
  \node[rung] (r3) at (0,-2.7)
    {\textbf{$+$ rewrites} \;$\longrightarrow$\; domain-specific languages};
  \draw[-{Latex[length=2mm]}] (r1) -- (r2);
  \draw[-{Latex[length=2mm]}] (r2) -- (r3);
\end{tikzpicture}
\end{center}

\section{Rung one: terms only}

Take $E = \emptyset$ and $R = \emptyset$. What remains is a multi-sorted signature: a
set of sorts and a set of operators with declared argument sorts and result sort. The
terms over such a signature, up to nothing at all, are exactly the inhabitants of an
algebraic data type.

\begin{lstlisting}
language! {
    name: Json,

    types {
        Value
        Field
    },

    terms {
        JNull  . Value ::= "null" ;
        JBool  . b:Bool |- b : Value ;
        JNum   . n:BigRat |- n : Value ;
        JStr   . s:Str |- s : Value ;
        JArr   . Value ::= "[" List(Value) "]" ;
        JObj   . Value ::= "{" List(Field) "}" ;
        Field  . k:Str, v:Value |- k ":" v : Field ;
    },

    equations { },
    rewrites  { },
}
\end{lstlisting}

There is nothing here a working developer has not written a hundred times, in whatever
notation their language provides for sums of products. The point of the rung is not
that it is novel but that it is \emph{included}: the specification format that will
shortly describe the $\pi$-calculus describes a JSON document with no change of
register.

\section{Rung two: equations}

Now populate $E$. The terms are quotiented by a congruence, and what one is presenting
is no longer a free term algebra but an algebra satisfying equations --- that is, a
finitely presentable object of universal algebra. Monoids, groups, rings, lattices,
semilattices, and the various bag- and set-like structures all live at this rung.

\begin{lstlisting}
language! {
    name: Monoid,

    types { M },

    terms {
        Unit . M ::= "e" ;
        Mul  . x:M, y:M |- x "*" y : M ;
    },

    equations {
        Assoc  . |- (Mul (Mul X Y) Z) = (Mul X (Mul Y Z)) ;
        UnitL  . |- (Mul Unit X)      = X ;
        UnitR  . |- (Mul X Unit)      = X ;
    },

    rewrites { },
}
\end{lstlisting}

The rung matters to us for a specific downstream reason. The equational theory of a
single operator --- whether it is free, associative, or associative--commutative ---
turns out to determine whether proximity to a resource confers authority over it
(Section~\ref{ob:subsec:capabilities}), and whether an interaction surface can be carried
by position or must be named explicitly (Section~\ref{ob:subsec:surfaces}). Equations are
not decoration. They are the difference between an object-capability discipline and an
ambient-authority leak.

\section{Rung three: rewrites}

Finally populate $R$. Now the data can wiggle.

This is the rung at which a specification stops describing a static structure and
starts describing a language with a behaviour. The reduction rules of an SOS
presentation go here, and so do the congruence rules that close reduction under the
term formers, in exactly the shape Plotkin, Milner, and Cardelli--Gordon wrote them.

\begin{lstlisting}
language! {
    name: Lambda,

    types { Term },

    terms {
        Lam . ^x.body:[Term -> Term] |- "lam " x "." body : Term;
        App . fun:Term, arg:Term     |- "(" fun "," arg ")" : Term;
    },

    equations { },

    rewrites {
        Beta     . |- (App (Lam fun) arg) ~> (eval fun arg);
        AppCongL . | M0 ~> M1 |- (App M0 N) ~> (App M1 N);
        AppCongR . | N0 ~> N1 |- (App M N0) ~> (App M N1);
        LamCong  . | S ~> T   |- (Lam ^x.S) ~> (Lam ^x.T);
    },
}
\end{lstlisting}

Read the \texttt{rewrites} block against a textbook presentation of the
$\lambda$-calculus and the correspondence is immediate: one base rule and three
congruence rules, with the binder handled by the \texttt{\^{}x.} notation and
substitution by \texttt{eval}. The horizontal bar of an inference rule has become
\texttt{|-}, and the premises sit to its left. This is not a translation of the
$\lambda$-calculus into some other formalism. It is the $\lambda$-calculus, written so
that a machine can read it.

\begin{remark}[What the ladder buys]
A reader arriving at rung three from rung one has crossed from data to computation
without changing tools. That is the whole design argument for the format, and it is
worth stating plainly: the reason to present a language as $(\Sigma, E, R)$ is not
elegance, it is that the same machinery that gives you a parser and a pretty-printer
for your JSON documents gives you a reduction engine for your calculus, and everything
in Part~III on top of that.
\end{remark}

\section{A worked presentation: the rho calculus}
\label{ob:subsec:rho-presentation}

We will need the rho calculus~\cite{meredith2005rho} as a running example, since
it is both the host language of Part~IV and one of the more instructive instances of
the constructions of Part~III. It is a reflective higher-order calculus: names are not
generated by a restriction operator but are \emph{quoted processes}, and the quoting and
unquoting operators $@$ and $*$ mediate between the two sorts.

\begin{lstlisting}
language! {
    name: RhoCalc,

    types { Proc  Name },

    terms {
        PZero   .                            |- "Nil" : Proc;
        PDrop   . n:Name                     |- "*" n : Proc;
        POutput . n:Name, q:Proc             |- n "!" "(" q ")" : Proc;
        PInput  . n:Name, ^x.p:[Name -> Proc] |- "for" "(" x "<-" n ")" p : Proc;
        PPar    . ps:HashBag(Proc)           |- "{" ps.*sep("|") "}" : Proc;
        NQuote  . p:Proc                     |- "@" p : Name;
    },

    equations {
        QuoteDrop . |- (NQuote (PDrop N)) = N ;
    },

    rewrites {
        Comm . |- (PPar {(PInput n ^x.p), (POutput n q), ...rest})
                    ~> (PPar {(subst ^x.p (NQuote q)), ...rest});
        Drop . |- (PDrop (NQuote P)) ~> P;
        ParCong . | S ~> T |- (PPar {S, ...rest}) ~> (PPar {T, ...rest});
    },
}
\end{lstlisting}

Three features of this presentation will be load-bearing later. The parallel composition
\texttt{PPar} takes a \texttt{HashBag}, so its equational theory is
associative--commutative --- position among parallel components carries no information.
The \texttt{Comm} rule matches two components of that bag and leaves the remainder
\texttt{...rest} untouched, which is what makes it a rule about a \emph{site} of
interaction rather than about whole terms. And the quote constructor \texttt{NQuote}
gives the theory a way to name its own programs, which will matter when we ask whether a
theory can internalise apparatus applied to it.

%===============================================================================
